% ABOUTME: Deep learning related work section of the IEEE report (Phase 2).
% ABOUTME: LSTM, TCN and Transformer, published FD001 deep results, protocol differences and design sources.
\section{Deep Learning Related Work}
\label{sec:dl-related}

Section~\ref{sec:related} placed our Phase 1 models among published FD001
results. Phase 2 trains networks that learn from raw sensor windows, so this
section reviews the three architectures we use, the deep models published on
C-MAPSS, what each architecture assumes about degradation data, and where
published numbers are scored differently from ours.

\subsection{Architectures}
\textbf{Recurrent networks and the LSTM.} A recurrent network updates a
hidden state once per cycle. Training it by backpropagation through time
multiplies the gradient by the recurrent weights once for every step it
travels back, so over long sequences the gradient shrinks towards zero or
grows without bound, and the network cannot learn dependencies more than a
few steps apart. Hochreiter and Schmidhuber~\cite{hochreiter1997lstm} fixed
this with a memory cell whose state is carried forward by addition rather
than by repeated multiplication. Multiplicative gates decide what is written
to the cell and what is read from it, and the forget gate added by Gers et
al.~\cite{gers2000forget} decides what is kept. While the forget gate stays
open, the error flowing back through the cell is passed on unchanged, so a
reading many cycles back can still affect the weights.

\textbf{Temporal convolutional networks.} Bai et al.~\cite{bai2018tcn}
proposed a generic convolutional architecture for sequences, the temporal
convolutional network (TCN). Its convolutions are causal: the output at
cycle $t$ uses only cycles up to $t$. They are also dilated: with dilation
$d$ and kernel size $k$, a filter at cycle $t$ reads cycles $t, t-d, \dots,
t-(k-1)d$, skipping the ones in between. Doubling $d$ from one residual block
to the next makes the receptive field grow exponentially with depth while
the number of weights grows only linearly. Across a range of sequence
modelling tasks, a generic TCN outperformed canonical LSTM and GRU networks
and kept a longer effective memory~\cite{bai2018tcn}.

\textbf{Transformers.} The Transformer of Vaswani et
al.~\cite{vaswani2017attention} replaces recurrence with self-attention.
Every position forms a query, a key and a value; its output is an average of
all values, weighted by how well its query matches each key. Any two cycles
are therefore one step apart, however far apart they are in time, and all
positions are processed in parallel. Self-attention ignores order, so
positional encodings are added to the inputs, and its cost grows with the
square of the window length.

\subsection{Deep models on C-MAPSS}
Babu et al.~\cite{babu2016cnn} applied the first convolutional network to
C-MAPSS. It reads 15-cycle windows of all sensors, chosen to fit the
shortest test history in their data, with a first filter spanning every
sensor and four cycles. On FD001 it reached RMSE 18.45 and score 1287.
Zheng et al.~\cite{zheng2017lstm} fed sensor sequences to stacked LSTM
layers with a training cap of 130~\cite{wang2019fda} and reached 16.14.
Li, Ding and Sun~\cite{li2018dcnn} used a one-dimensional CNN on 30-cycle
windows of the 14 informative sensors with a cap of 125: four convolution
layers of ten filters of length 10, one of length 3 and a fully connected
layer of 100 units. Its RMSE of 12.61 is the lowest in the comparison of Ragab et
al.~\cite{ragab2020attention}.

Attention entered with the encoder-decoder ATS2S of Ragab et
al.~\cite{ragab2020attention}, which attends over the hidden states of an
LSTM encoder (RMSE 12.63, score 243). Mo et al.~\cite{mo2021gcu} use a
Transformer encoder as the backbone and add a gated convolutional unit,
arguing that attention alone makes the learned features insensitive to local
context; later comparisons report an FD001 RMSE of 11.27 for this model.
Zhang, Song and Li~\cite{zhang2022dast} propose DAST, built only from
self-attention: two encoders run in parallel, one attending across sensors
and one across time steps. With 40-cycle windows and a cap of 125 it reports
RMSE 11.43 and score 203.15 on FD001.

Hybrid models combine learned and engineered features. Chen et
al.~\cite{chen2021attention} weight the hidden states of an LSTM over
30-cycle windows with an attention layer and concatenate the result with two
handcrafted features of each window, its mean and the coefficient of a
linear trend; they report RMSE 14.53 and score 322.44 on FD001, as tabulated
in~\cite{zhang2022dast}. Liu et al.~\cite{liu2021rul} instead pass features
learned by a deep CNN to a LightGBM regressor. The Phase 2 hybrid
model follows the first route, feeding our Phase 1 features to a sequence
encoder.

Table~\ref{tab:dl-published} collects these results. The full table, with
the source of each value, is \texttt{results/dl\_published.csv}.

\begin{table}[t]
\centering
\caption{Published FD001 deep learning results. Values for the first seven
rows are those collected in~\cite{ragab2020attention}, Table III; Chen et
al.\ and DAST are from~\cite{zhang2022dast}. $w$ is the window length and
Cap the training RUL cap. A dash marks a value not stated in the sources we
checked. None of the sources we checked states whether test labels were
clipped before scoring. Our models are scored against labels clipped at 125
(Table~\ref{tab:published}).}
\label{tab:dl-published}
\footnotesize
\begin{tabular}{lrrrr}
\hline
Method & $w$ & Cap & RMSE & $S$ \\
\hline
2D CNN \cite{babu2016cnn}           & 15 & --  & 18.45 & 1287 \\
Deep LSTM \cite{zheng2017lstm}      & -- & 130 & 16.14 & 338 \\
MODBNE \cite{zhang2017modbne}       & -- & --  & 15.04 & 334 \\
CNN-LSTM, Wu et al.                 & -- & --  & 14.40 & 290 \\
HDNN, Al-Dulaimi et al.             & -- & --  & 13.02 & 245 \\
ATS2S \cite{ragab2020attention}     & 30 & 125 & 12.63 & 243 \\
1D CNN \cite{li2018dcnn}            & 30 & 125 & 12.61 & 274 \\
\hline
LSTM-att.\ + features \cite{chen2021attention} & 30 & -- & 14.53 & 322 \\
GCU-Transformer \cite{mo2021gcu}    & -- & --  & 11.27 & -- \\
DAST \cite{zhang2022dast}           & 40 & 125 & 11.43 & 203 \\
\hline
Random Forest (ours, Phase 1)       & -- & 125 & 14.02 & 310 \\
XGBoost (ours, Phase 1)             & -- & 125 & 14.08 & 347 \\
\hline
\end{tabular}
\end{table}

\subsection{What each architecture assumes}
An FD001 sensor trace is flat and noisy while the engine is healthy, then
drifts steadily as the compressor degrades. Each architecture brings a
different assumption about how to read that curve. The LSTM assumes the
useful history can be summarised in a fixed-size state updated one cycle at
a time. That suits damage that accumulates, but whatever the network learns
about early cycles must survive every later update, so recent cycles tend to
dominate. A CNN or TCN assumes the useful patterns are local and look the
same wherever they occur: one filter detects a rise or a step anywhere in
the window. That suits noisy sensors, where smoothing and a local slope carry
most of the signal, and it is what our Phase 1 rolling means and slopes do
by hand; but nothing outside the receptive field can affect the output, so
a drift longer than the receptive field is invisible. Self-attention assumes
nothing about locality or order: any cycle may matter to any other, and
the weights are computed from the readings themselves, so the model can
focus on the cycles where degradation begins. The price is that the
structure must be learned from data, and FD001 has only 100 training
engines. Mo et al.\ and DAST both add structure back, through a local
convolution or a separate attention over sensors~\cite{mo2021gcu,
zhang2022dast}, which suggests that plain attention underuses the local
smoothness of degradation curves.

\subsection{How published numbers are scored}
Published FD001 numbers are not all produced under one protocol, and the
differences are larger than many of the gaps between models.
\begin{itemize}
\item \emph{Test labels.} Clipping the test RUL at the cap moves our
XGBoost RMSE from 15.36 to 14.08. None of the sources we checked states
which labels it scores against.
\item \emph{Cap and window.} Caps of 125 and 130 and windows of 15 to 40
cycles are all in use. The shortest FD001 test engine has 31 cycles, so a
40-cycle window needs padding or must drop engines, which changes the set
being scored; how this is handled is not always reported.
\item \emph{Selection and repeats.} Some papers report the best of several
runs, or do not state how hyperparameters were chosen. Ours are chosen by
cross validation on training engines only, and the test set is scored once.
\item \emph{Score $S$.} $S$ is a sum over engines, so it is only comparable
on the same 100 test engines.
\item \emph{Transcribed values.} We copy published numbers, often from
another paper's comparison table, rather than rerun the models.
\end{itemize}
Gaps below about one cycle of RMSE should therefore not be read as
meaningful.

\subsection{Sources for our design choices}
The sequence models of Sections~\ref{sec:dl-setup} and~\ref{sec:dl-models}
follow published practice. The 30-cycle window is that of Li et
al.~\cite{li2018dcnn}, ATS2S~\cite{ragab2020attention} and Chen et
al.~\cite{chen2021attention}, and fits the shortest FD001 test engine. The
cap of 125 follows the piecewise-linear target of
Heimes~\cite{heimes2008rnn} with the value used for FD001
in~\cite{li2018dcnn,ragab2020attention,zhang2022dast}. Dropout follows
Srivastava et al.~\cite{srivastava2014dropout}; the dilated causal residual
blocks of the TCN follow Bai et al.~\cite{bai2018tcn}; the LSTM uses the
forget gate of Gers et al.~\cite{gers2000forget}; and the Transformer's
sinusoidal positional encoding is that of Vaswani et
al.~\cite{vaswani2017attention}.
